\documentclass[10pt,a4paper]{amsart}
\usepackage[margin=19mm]{geometry}
\usepackage{amsmath,amssymb,mathtools}
\usepackage[scr=rsfs,cal=euler]{mathalfa}
\usepackage{xfrac}
\usepackage[unicode,hidelinks,bookmarks=false]{hyperref}
\newcommand{\ie}{i.e.\ }
\newcommand{\cf}{cf.\ }
\newcommand{\resp}{respectively\ }
\newcommand{\Subsection}{\subsection}
\providecommand{\nobreakdash}{\nobreak\hbox{-}}
\setlength{\emergencystretch}{2em}
\multlinegap0pt
\usepackage{xfrac}
\usepackage{dsfont}
\usepackage{xcolor}
\usepackage{amssymb}
\usepackage{float}
\usepackage[shortlabels]{enumitem}
\newtheorem{thm}{Theorem}
\newtheorem{lem}{Lemma}
\newtheorem{cor}{Corollary}
\newtheorem{rem}{Remark}
\newtheorem{defn}{Definition}
\DeclareMathOperator{\Ld}{\mathcal{L}^d}
\newcommand{\bftheta}{{\boldsymbol\theta}}
\newcommand{\diag}{\mathop{\mathrm{diag}}}\newcommand{\esssup}{\mathop{\mathrm{ess\,sup}}}\newcommand{\essinf}{\mathop{\mathrm{ess\,inf}}}\newcommand{\e}{\textup{e}}\renewcommand{\i}{\textup{i}}\newcommand{\bfa}{\mathbf a}\newcommand{\bfb}{\mathbf b}\newcommand{\bfc}{\mathbf c}\newcommand{\bfd}{\mathbf d}\newcommand{\bfe}{\mathbf e}\newcommand{\bff}{\mathbf f}\newcommand{\bfg}{\mathbf g}\newcommand{\bfh}{\mathbf h}\newcommand{\bfi}{\mathbf i}\newcommand{\bfj}{\mathbf j}\newcommand{\bfk}{\mathbf k}\newcommand{\bfl}{\mathbf l}\newcommand{\bfm}{\mathbf m}\newcommand{\bfn}{\mathbf n}\newcommand{\bfo}{\mathbf o}\newcommand{\bfp}{\mathbf p}\newcommand{\bfq}{\mathbf q}\newcommand{\bfr}{\mathbf r}\newcommand{\bfs}{\mathbf s}\newcommand{\bft}{\mathbf t}\newcommand{\bfu}{\mathbf u}\newcommand{\bfw}{\mathbf w}\newcommand{\bfv}{\mathbf v}\newcommand{\bfx}{\mathbf x}\newcommand{\bfy}{\mathbf y}\newcommand{\bfz}{\mathbf z}\newcommand{\bfzero}{\mathbf 0}\newcommand{\bfuno}{\mathbf 1}
\title{The $*$-algebra of unbounded GLT: construction and theoretical foundations}
\author{Andrea Adriani, Alec Jacopo Almo Schiavoni-Piazza}
\date{Source: arXiv:2602.23879v1 (CC BY 4.0)}
\begin{document}
\maketitle

\noindent\emph{Source and licence.} The $*$-algebra of unbounded GLT: construction and theoretical foundations. Version arXiv:2602.23879v1 (CC BY 4.0), distributed by arXiv under the Creative Commons Attribution 4.0 International licence, http://creativecommons.org/licenses/by/4.0/, as recorded in the arXiv licence metadata for this exact version and shown on the abstract page https://arxiv.org/abs/2602.23879v1. Full licence notice: \texttt{licenses/arxiv-2602.23879-CC-BY-4.0.txt}.\par\smallskip

\noindent\textit{Editorial scope.} Counted unit: the article's thm thm\_acs\_closure (source0.tex lines 1628--1635). The article's complete original proof of it is delivered with it (source0.tex lines 1636--1785, one \texttt{proof} environment of 150 lines). No mathematics is written, repaired or re-derived: the statement and the proof are the article's own lines, in the article's own notation, and no retained line is substituted or re-flowed.  Retained uncounted as the source-local context the proof needs: lines 228--238; lines 488--508; lines 532--542; lines 565--567; lines 568--570; lines 593--599; lines 679--698; lines 814--823; lines 827--845; lines 860--874; lines 966--968; lines 1036--1073; lines 1109--1134; lines 1222--1236; lines 1345--1354; lines 1470--1477.  Nothing was omitted from any retained span, so the package carries no omission record.  No retained line contains a citation, so the wrapper carries no bibliography and no bibliography entry.  No retained line contains a figure environment, an image-inclusion command or drawing code, so no image is delivered and the package declares no asset dependency.  Editorial formatting only, in addition: an AMS \texttt{amsart} preamble that restates the article's own declarations and macros used by the retained lines, namely the environments \texttt{thm}, \texttt{lem}, \texttt{cor}, \texttt{rem}, \texttt{defn} on separate counters, together with the standard packages those lines need.  The retained environments print in this document as Theorem 1, Lemma 1, Corollary 1, Remark 1, Definition 1 in order of appearance, whereas the article prints the same items with the numbering of its own full document.  The complete native source of the article as distributed in the arXiv e-print for this exact version is bundled unchanged as \texttt{sources/195429/source0.tex}.
\begin{thm}\label{completeness_acs_conv}
    The space $(\mathcal{M},d_{a.c.s.})$ is a complete pseudo-metric space and, given $\{A_n\}_n,\{B_n\}_n\in\mathcal{M}$, it holds
    \begin{equation*}
        d_{a.c.s.}(\{A_n\}_n,\{B_n\}_n)=0\quad\iff\quad\{A_n-B_n\}_n\sim_{\sigma} 0.
    \end{equation*}
    Moreover, given $\{A_n\}_n,\{\{B_{n,t}\}_n\}_t\in\mathcal{M}$, the following are equivalent:
    \begin{itemize}
        \item $\lim\limits_{t\to+\infty} d_{a.c.s.}(\{B_{n,t}\}_n,\{A_n\}_n)=0;$
        \item $\{\{B_{n,t}\}_n\}_t$ is an a.c.s. for $\{A_n\}_n$.
    \end{itemize}
\end{thm}

\begin{lem}\label{lemma_1}
	 Let $Q_1=Q_{{\bfy_1},l_1}\subset Q_2=Q_{{\bfy_2},l_2}$ be two hypercubes. For any multi-index $\bfn$, we have 
	\begin{equation}\label{lemma_1_eq_1}
	 	\Pi_{\bfn}\Pi_{\bfn}^T = I_{\bfn}\in\mathbb{C}^{N(l_1\cdot \bfn)\times N(l_1\cdot\bfn)},
	\end{equation}
	\begin{equation}\label{lemma_1_eq_2}
		\Pi_{\bfn}^T \Pi_{\bfn} = D_{l_2\bfn}(\mathds{1}_{Q_1})\in\mathbb{C}^{N(l_2\cdot \bfn)\times N(l_2\cdot\bfn)},
	\end{equation}
	where $I_{\bfn}$ is the identity matrix of size $N(l_1\bfn)$ and $D_{l_2\bfn}(\mathds{1}_{Q_1})$ is the $l_2\bfn$-th matrix of the diagonal matrix-sequence $\{D_{l_2\bfn}(\mathds{1}_{Q_1})\}_\bfn\in\mathcal{G}_{{\bfy_2},l_2}$, generated by the characteristic function of $Q_1$, namely
	\begin{equation*}
		D_{l_2\bfn}(\mathds{1}_{Q_1})=\diag\biggl(\mathds{1}_{Q_1}\biggl(\bfy_2+\frac{\bfi}{\bfn}\biggr)\biggr)_{{\bfi}={\bf1},\ldots,l_2\cdot\bfn}.
	\end{equation*}
	As a consequence, for square matrices $A_{\bfn}$ of size $N(l_2\bfn)$ and $B_{\bfn}$ of size $N(l_1\bfn)$, we have:
	\begin{equation*}
		R_{\bfn}(E_{\bfn}(B_{\bfn}))=B_{\bfn},
	\end{equation*}
	\begin{equation*}
		E_{\bfn}(R_{\bfn}(A_{\bfn}))=D_{l_2\bfn}(\mathds{1}_{Q_1})A_{\bfn}D_{l_2\bfn}(\mathds{1}_{Q_1}).
	\end{equation*}

\end{lem}

\begin{cor}\label{product_hypercubes}
     Let $Q_1=Q_{{\bfy_1},l_1}\subset Q_2=Q_{{\bfy_2},l_2}$ be two hypercubes and let $\bfn$ be a multi-index. Then, if $B_\bfn,B'_\bfn$ are square matrices of size $N(l_1\bfn)$, it holds
     \begin{equation*}
         E_\bfn(B_\bfn B'_\bfn)=E_\bfn(B_\bfn)E_\bfn(B'_\bfn).
     \end{equation*}
     Conversely, if $A_\bfn,A'_\bfn\in\mathcal{C}^{N(l_2\cdot\bfn)\times N(l_2\cdot\bfn)}$ are in the image of the extension operator $E_\bfn$, then
     \begin{equation*}
         R_\bfn(A_\bfn A'_\bfn)=R_\bfn(A_\bfn)R_\bfn(A'_\bfn).
     \end{equation*}
     As a consequence, the extension operator $E_\bfn$ is a $*$-algebra morphism, whilst the restriction operator $R_\bfn$ is a $*$-algebra morphism, when restricted to the image of the extension operator. 
\end{cor}

\begin{lem}\label{acs_cts_operators}
    Let $Q_1(=Q_{{\bfy_1},l_1})\subset Q_2(=Q_{{\bfy_2},l_2})$ be two hypercubes. Denote by $\mathcal{G}_1=\mathcal{G}_{{\bfy_1},l_1}$ and $\mathcal{G}_2=\mathcal{G}_{{\bfy_2},l_2}$ the corresponding GLT algebras and by $\mathcal{M}_1$ and $\mathcal{M}_2$ the algebras of all matrix-sequences of size $\{l_1^d\cdot N(\bfn)\}_\bfn$ and $\{l_2^d\cdot N(\bfn)\}_\bfn$, respectively. Then, the operators $\mathcal{R}=\{R_{\bfn}\}_\bfn:\mathcal{M}_2\to\mathcal{M}_1$ and $\mathcal{E}=\{E_{\bfn}\}_\bfn:\mathcal{M}_1\to\mathcal{M}_2$ are continuous with respect to a.c.s. convergence.
\end{lem}

\begin{rem}\label{acs_same_structure}
    In the following sections, we make use of many other restriction and extension operators, built with the same structure as those between two hypercubes. Therefore, all those operators will be continuous with respect to a.c.s. convergence, with the same proof as the one of Lemma \ref{acs_cts_operators}, as long as the sequences of dimensions have the same rate of divergence (i.e. their ratio has a positive finite limit). 
\end{rem}

\begin{thm}\label{restriction_of_standard_GLT}
	Let $Q_1(=Q_{{\bfy_1},l_1})\subset Q_2(=Q_{{\bfy_2},l_2})$ be two hypercubes.
    If $\mathcal{R}=\{R_{\bfn}\}_\bfn:\mathcal{M}_2\to\mathcal{M}_1$ and $\mathcal{E}=\{E_{\bfn}\}_\bfn:\mathcal{M}_1\to\mathcal{M}_2$ are the restriction and extension operators, then we have
	\begin{equation}\label{GLT_op}
		\mathcal{R}(\mathcal{G}_2)\subset\mathcal{G}_1, \qquad \mathcal{E}(\mathcal{G}_1)\subset\mathcal{G}_2.
	\end{equation}
\end{thm}

\begin{rem}\label{rest_surjective}
    Note that, given two hypercubes $Q_1\subset Q_2$, by Lemma \ref{lemma_1}, we have
    \begin{equation*}
        \mathcal{R}\circ\mathcal{E}=\mathrm{Id}_{\mathcal{M}_1},
    \end{equation*}
    where $\mathrm{Id}_{\mathcal{M}_1}:\mathcal{M}_1\to\mathcal{M}_1$ is the identity operator. In particular, this implies that
    \begin{equation*}
        (\mathcal{R}\circ\mathcal{E})(\mathcal{G}_1)=\mathcal{G}_1.
    \end{equation*}
    As a consequence, we can update Theorem \ref{restriction_of_standard_GLT} obtaining that
    \begin{equation*}
        \mathcal{R}(\mathcal{G}_2)=\mathcal{G}_1,
    \end{equation*}
    i.e., the restriction operator is always surjective. 
    
    Similarly, one can easily prove that the image of the extension operator is the sub-algebra of $\mathcal{G}_2$ of sequences whose canonical symbol $f$ satisfies
    \begin{equation*}
        f(\bfx,\bftheta)=0,\qquad \text{for a.e. }\bfx\in Q_2\setminus Q_1.
    \end{equation*}
    \end{rem}

\begin{defn}[Reduced GLT algebra]\label{reduced_GLT}
    Consider an open bounded domain $\Omega$ such that $\mathcal{L}^d(\partial\Omega)=0$ and fix a hypercube $Q_{\bfy,l}$ such that $\Omega\subset Q_{\bfy,l}$. 
    
    Let $\mathcal{M}_{\bfy,l}$ and $\mathcal{M}_{\Omega}$ be the sets of all square matrix-sequences of size $\{l^d\cdot N(\bfn)\}_{\bfn}$ and $\{d_{\bfn}^{\Omega}\}_{\bfn}$, respectively. 
    
    If $\mathcal{R}_{\bfy,l,\Omega}=\{R_\bfn\}_\bfn:\mathcal{M}_{\bfy,l}\to\mathcal{M}_{\Omega}$ is the restriction operator, we define the reduced GLT algebra on $\Omega$, and we denote it by $\mathcal{G}_{\Omega}$, as
    \begin{equation*}
        \mathcal{G}_{\Omega}:=\mathcal{R}_{\bfy,l,\Omega}(\mathcal{G}_{\bfy,l}).
    \end{equation*}
\end{defn}

\begin{rem}\label{well-posedness}
    Note that, thanks to Theorem \ref{restriction_of_standard_GLT}, the definition of reduced GLT algebra over a regular bounded domain $\Omega$ does not depend on the choice of the hypercube $Q_{\bfy,l}$ containing $\Omega$. Indeed, thanks to the compatibility between the two definitions of restriction operators (between two hypercubes and from a hypercube to any bounded open domain), they factor through GLT algebras built upon intermediate hypercubes. 
    
    More precisely, if $\Omega\subset Q_{\bfy_1,l_1}\subset Q_{\bfy_2,l_2}$ and 
    \begin{align*}
        \mathcal{R}_{\bfy_2,l_2,\bfy_1,l_1}&:\mathcal{M}_{\bfy_2,l_2}\to\mathcal{M}_{\bfy_1,l_1},\\
        \mathcal{R}_{\bfy_2,l_2,\Omega}&:\mathcal{M}_{\bfy_2,l_1}\to\mathcal{M}_{\Omega},\\
        \mathcal{R}_{\bfy_1,l_1,\Omega}&:\mathcal{M}_{\bfy_1,l_1}\to\mathcal{M}_{\Omega}
    \end{align*}
    are the restriction operators as defined above, then 
    \begin{equation*}
         \mathcal{R}_{\bfy_1,l_1,\Omega}\circ\mathcal{R}_{\bfy_2,l_2,\bfy_1,l_1}=\mathcal{R}_{\bfy_2,l_2,\Omega}
    \end{equation*}
    and 
    \begin{equation*}
        \mathcal{R}_{\bfy_2,l_2,\bfy_1,l_1}(\mathcal{G}_{\bfy_2,l_2})=\mathcal{G}_{\bfy_1,l_1},
    \end{equation*}
    so that the image of the two operators is the same. This suggests that Definition \ref{reduced_GLT} selects matrix-sequences whose spectral properties are related to the domain $\Omega$, which is the case, as proven in Theorem \ref{spectra_reduced_GLT}. 
\end{rem}

\begin{thm}\label{spectra_reduced_GLT}
    Consider an open bounded domain $\Omega$, with $\mathcal{L}^d(\partial\Omega)=0$ and a hypercube $Q_{\bfy,l}$ such that $\Omega\subset Q_{\bfy,l}$. 
    
    For any $\{B_{\bfn}\}_\bfn\in\mathcal{G}_{\Omega}$, consider a sequence $\{A_{\bfn}\}_\bfn$ in its preimage through the restriction map $\mathcal{R}:\mathcal{G}_{\bfy,l}\to\mathcal{G}_{\Omega}$, with canonical GLT symbol $f:Q_{\bfy,l}\times[-\pi,\pi]^d\to\mathbb{C}$.\\
    Then
    \begin{equation*}
        \{B_{\bfn}\}_{\bfn}\sim_{\sigma}\left(g,\Omega\times\left[-\pi,\pi\right]^d\right),
    \end{equation*}
    with $g=f|_{\Omega\times[-\pi,\pi]^d}$.
    
    If, additionally, $\{B_\bfn\}_\bfn$ is a sequence of Hermitian matrices, then it also holds 
    \begin{equation*}
        \{B_{\bfn}\}_{\bfn}\sim_{\lambda}\left(g,\Omega\times\left[-\pi,\pi\right]^d\right).
    \end{equation*}
\end{thm}

\begin{thm}\label{isometry_red_GLT}
 Let $\Omega$ be an open bounded domain such that $\Ld(\partial\Omega)=0$, and consider the map $\Phi_\Omega:(\sfrac{\mathcal{G}_\Omega}{\sim_d},d_{a.c.s.})\to(L^0(\Omega\times[-\pi,\pi]^d),d_{m,\Omega})$ that associates to any reduced GLT its canonical symbol. Then, $\Phi_\Omega$ is a surjective isometry.   
\end{thm}

\begin{rem}\label{rem_restr_btw_domains}
    If we consider an open domain $\Omega$, such that $d_\bfn^{\Omega}=\#\Theta_{\bfn,\Omega}<+\infty$, we can naturally define the algebra of all matrix-sequences of size $\{d_\bfn^{\Omega}\}_\bfn$.
    
    In addition, if two domains $\Omega_1\subset\Omega_2$ are such that $d_\bfn^{\Omega_i}=\#\Theta_{\bfn,\Omega_i}<+\infty$ for $i=1,2$ and for every $\bfn$, then we can define restriction and extension operators between $\mathcal{M}_{\Omega_1}$ and $\mathcal{M}_{\Omega_2}$. 
    
    Indeed, $\Omega_1\subset\Omega_2$ implies that $\Theta_{\bfn,\Omega_1}\subset\Theta_{\bfn,\Omega_2}$. Thus, consider the family of matrices $\Pi_{\bfn}(=\Pi_{\bfn,\Omega_2,\Omega_1})\in\mathbb{C}^{d_{\bfn}^{\Omega_1}\times d_{\bfn}^{\Omega_2}}$, defined by
\begin{equation*}
    (\Pi_{\bfn})_{\bfi,\bfj}=\delta_{\bfi,\bfj},\qquad\bfi\in\Theta_{\bfn,\Omega_1},\quad \bfj\in\Theta_{\bfn,\Omega_2}.
\end{equation*}
Similarly as above, the restriction operator $\mathcal{R}_{\Omega_2,\Omega_1}=\{R_{\bfn}\}_\bfn:\mathcal{M}_{\Omega_2}\to\mathcal{M}_{\Omega_1}$ and the extension operator $\mathcal{E}_{\Omega_1,\Omega_2}=\{E_{\bfn}\}_\bfn:\mathcal{M}_{\Omega_1}\to\mathcal{M}_{\Omega_2}$ are defined as 
\begin{equation*}
	R_{\bfn}(=R_{\bfn,\Omega_2,\Omega_1}):\mathbb{C}^{d_{\bfn}^{\Omega_2}\times d_{\bfn}^{\Omega_2}}\to\mathbb{C}^{d_{\bfn}^{\Omega_1}\times d_{\bfn}^{\Omega_1}},
\end{equation*}
\begin{equation*}
	E_{\bfn}(=E_{\bfn,\Omega_1,\Omega_2}):\mathbb{C}^{d_{\bfn}^{\Omega_1}\times d_{\bfn}^{\Omega_1}}\to\mathbb{C}^{d_{\bfn}^{\Omega_2}\times d_{\bfn}^{\Omega_2}},
\end{equation*}
where, explicitly,
\begin{equation*}
	R_{\bfn}(A_{\bfn}):=\Pi_{\bfn}A_{\bfn}\Pi_{\bfn}^T,
\end{equation*}
\begin{equation*}
	E_{\bfn}(B_{\bfn}):=\Pi_{\bfn}^T B_{\bfn}\Pi_{\bfn}. 
\end{equation*}

The operators $\mathcal{R}_{\Omega_2,\Omega_1}$ and $\mathcal{E}_{\Omega_1,\Omega_2}$ share the same structure of restriction and extension operators between a hypercube $Q$ and a domain $\Omega$. For this reason, all operators factor through ``intermediate domains'' as in Remark \ref{well-posedness}.

Indeed, this compatibility observation yields
\begin{align*}
    \mathcal{R}_{\Omega_2,\Omega_1}(\mathcal{G}_{\Omega_2})&\subset\mathcal{G}_{\Omega_1},\\
    \mathcal{E}_{\Omega_1,\Omega_2}(\mathcal{G}_{\Omega_1})&\subset\mathcal{G}_{\Omega_2}.
\end{align*}
Moreover, reasoning similarly as in Remark \ref{rest_surjective}, one can prove that the restriction operator is surjective, whilst the image of the extension operator is the sub-algebra of $\mathcal{G}_{\Omega_2}$ of sequences whose canonical symbol satisfies 
\begin{equation*}
        f(\bfx,\bftheta)=0,\qquad \text{for a.e. }\quad(\bfx,\bftheta)\in (\Omega_2\setminus \Omega_1)\times[-\pi,\pi]^d.
    \end{equation*}
For the same reason as in Corollary \ref{product_hypercubes}, the extension operator $\mathcal{E}_{\Omega_1,\Omega_2}$ is a $*$-algebra morphism, whilst the restriction operator $\mathcal{R}_{\Omega_2,\Omega_1}$ is a $*$-algebra morphism, when restricted to the image of $\mathcal{E}_{\Omega_1,\Omega_2}$.
Finally, by Lemma \ref{acs_cts_operators} and Remark \ref{acs_same_structure}, both the extension and the restriction operators are continuous with respect to the a.c.s. convergence, also in this case, as long as the sequences of dimensions have the same rate of divergence (note that for the class of domains of interest, this will always be the case, as a byproduct of item $(iii)$ of Lemma \ref{dimensions}).
\end{rem}

\begin{rem}\label{rem_obs_ext_operator}
    In the next section, we will need the following useful observation.
    Consider two domains $\Omega_1\subset\Omega_2$, with $\Ld(\partial\Omega_i)=0$, for $i=1,2$. Fix any $\bfn$ and any $B_\bfn\in\mathcal{C}^{d_\bfn^{\Omega_1}\times d_\bfn^{\Omega_1}}$ and consider the extension operator $E_{\bfn}:\mathbb{C}^{d_{\bfn}^{\Omega_1}\times d_{\bfn}^{\Omega_1}}\to\mathbb{C}^{d_{\bfn}^{\Omega_2}\times d_{\bfn}^{\Omega_2}}$. 
    
    Define $\tilde{\Pi}_\bfn\in\mathbb{C}^{d_{\bfn}^{\Omega_2}\times d_{\bfn}^{\Omega_2}}$ as 
    \begin{equation*}
           \tilde{\Pi}_\bfn:= \left(
    \begin{array}{c} 
      \Pi_\bfn \\
      \hline
      C_\bfn
    \end{array} 
    \right), 
    \end{equation*}
    where $C_\bfn(=C_{\bfn,\Omega_2,\Omega_1})$ is any matrix of size $(d_\bfn^{\Omega_2}-d_\bfn^{\Omega_1})\times d_\bfn^{\Omega_2}$ that completes $\Pi_\bfn$ to a permutation matrix (for any row of $C_\bfn$, set all but one entry to $0$, and the last entry equal to $1$, in such a way that any column of $\tilde{\Pi}_\bfn$ has just one non-zero entry). Then, 
    \begin{equation*}
        E_\bfn(B_\bfn)=\Pi_\bfn^{T}B_\bfn\Pi_\bfn=\tilde{\Pi}_\bfn^{T}\tilde{B}_\bfn\tilde{\Pi}_\bfn,
    \end{equation*}
    with $\tilde{B}_\bfn:=B_\bfn\oplus0_\bfn$ and $0_\bfn(=0_{\bfn,\Omega_2,\Omega_1})$ is a square null matrix of size $d_\bfn^{\Omega_2}-d_\bfn^{\Omega_1}$.
    
    As a consequence, any square matrix $A_\bfn$ of size $d_\bfn^{\Omega}$ can be written as 
    \begin{equation}\label{restr_ext_identity}
        A_\bfn=E_\bfn(R_\bfn(A_\bfn))+S_\bfn=\tilde{\Pi}_\bfn^{T}\left(R_\bfn(A_\bfn)\oplus0_\bfn\right)\tilde{\Pi}_\bfn+S_\bfn,
    \end{equation}
    with $S_\bfn(=S_{\bfn,\Omega_2,\Omega_1})$ (which in the following will play the role of a rank correction) of rank bounded by $2(d_\bfn^{\Omega_2}-d_\bfn^{\Omega_1})$, since it has at most $d_\bfn^{\Omega_2}-d_\bfn^{\Omega_1}$ rows and columns with non-zero entries.
\end{rem}

\begin{lem}\label{dimensions}
    Let $\Omega$ be an open domain, such that $\Ld(\Omega)<\infty$, $\Ld(\partial\Omega)=0$. Consider a regular exhaustion $\{\Omega_t\}_t$ of $\Omega$.\\
    \begin{itemize}
        \item[(i)]\label{dim_item_1} There exist $\bfn_0$ and $c>0$, such that $d_\bfn^{\Omega}\geq cN(\bfn)$, for every $\bfn\geq\bfn_0$;\\
        \item[(ii)]\label{dim_item_2} there exists $c'>0$ such that, for every $t$ and every $\bfn\geq\bfn_t$, it holds 
        \begin{equation*}
            \frac{d_\bfn^{\Omega}-d_\bfn^{\Omega_t}}{d_\bfn^{\Omega}}\leq c'\Ld(\Omega\setminus\Omega_t);
        \end{equation*}
        \item[(iii)]\label{dim_item_3} as $\bfn\to\infty$, it holds
        \begin{equation*}
            \frac{d_\bfn^{\Omega}}{N(\bfn)}\to\Ld(\Omega).
        \end{equation*}
    \end{itemize}
    %\textcolor{blue}{non so se mi servono tutti, ma intanto li dimostro e li abbiamo. Inolre, non mi piace troppissimo la spaziatura di itemize e enumerate in sto template, possiamo farci qualcosa? Confermo che l'item (iii) non credo servirà nel seguito, cmq per ora lo terrei e anzi può essere interessante di per se, forse.}
\end{lem}

\begin{defn}[Unbounded GLT sequences]\label{def_UGLT}
    Let $\Omega$ be an open domain, such that $\Ld(\Omega)<\infty$, $\Ld(\partial\Omega)=0$.
    Consider the algebra $\mathcal{M}_{\Omega}$ of all matrix-sequences of size $\{d_\bfn^{\Omega}\}_\bfn$. An element $\{A_\bfn\}_\bfn\in\mathcal{M}_\Omega$ is called an unbounded GLT over $\Omega$ if there exists a regular exhaustion $\{\Omega_t\}_t$ of $\Omega$ and a sequence of matrix-sequences $\left\{\{B_{\bfn,t}\}_\bfn\right\}_t$ of size $\left\{\{d_\bfn^{\Omega_t}\}_\bfn\right\}_t$, such that
    \begin{itemize}
        \item[(i)] $\forall t$, $\{B_{\bfn,t}\}_\bfn\in\mathcal{G}_{\Omega_t}$, with GLT symbol $f_t:\Omega_t\times[-\pi,\pi]^d\to\mathbb{C}$;
        \item[(ii)] $\left\{\{B_{\bfn,t}\}_\bfn\right\}_t$ is a $\mathrm{g.a.c.s.}$ for $\{A_\bfn\}_\bfn$, with structure matrix-sequences given by $\left\{\left\{\tilde{\Pi}^{T}_{\bfn,t}\right\}_\bfn\right\}_t$ and $\left\{\left\{\tilde{\Pi}_{\bfn,t}\right\}_\bfn\right\}_t$, as in Remark \ref{rem_obs_ext_operator};
        \item[(iii)] $\forall t<t'$, $f_t=f_{t'}$, almost everywhere in  $\Omega_t\times[-\pi,\pi]^d$.
    \end{itemize}
    The sequence $\{\{B_{\bfn,t}\}_\bfn\}_t$ is regarded as an approximating sequence for $\{A_\bfn\}_\bfn$, associated with the exhaustion $\{\Omega_t\}_t$.
\end{defn}

\begin{thm}[Permanence of spectral information]\label{perm_spectr_info}
    Let $\Omega$ be an open domain, such that $\Ld(\Omega)<\infty$, $\Ld(\partial\Omega)=0$. Consider an unbounded GLT $\{A_\bfn\}_\bfn\in\mathcal{M}_\Omega$. Then, for every open bounded domain $\Omega_0\subset\Omega$ such that $\Ld(\partial\Omega_0)=0$, it holds $\mathcal{R}_{\Omega,\Omega_0}(\{A_\bfn\}_\bfn)\in\mathcal{G}_{\Omega_0}$.
    
    Additionally, let $f:\Omega\times[-\pi,\pi]^d\to\mathbb{C}$ be any symbol for $\{A_\bfn\}_\bfn$, computed via an approximating sequence defining $\{A_\bfn\}_\bfn$. Then
    \begin{equation*}
        \mathcal{R}_{\Omega,\Omega_0}(\{A_\bfn\}_\bfn)\sim_{GLT}^{\Omega_0}f|_{\Omega_0\times[-\pi,\pi]^d}.
    \end{equation*}
\end{thm}

\begin{thm}[a.c.s. closure]\label{thm_acs_closure}
    Let $\Omega$ be an open domain, such that $\Ld(\Omega)<+\infty$ and $\Ld(\partial\Omega)=0$. Consider a sequence of unbounded $\mathrm{GLT}$ $\{\{A_{\bfn,s}\}_\bfn\}_s\sim_{\mathrm{GLT}}^{\Omega}f_s$, over $\Omega$. Then, the following are equivalent:
    \begin{itemize}
        \item[(i)]\label{thm_acs_closure_item_i} There exists $\{A_\bfn\}_\bfn\in\mathcal{M}_\Omega$, such that $\{\{A_{\bfn,s}\}_\bfn\}_s\overset{a.c.s.}{\longrightarrow}\{A_\bfn\}_\bfn$; 
        \item[(ii)]\label{thm_acs_closure_item_ii} There exists $f\in\mathcal{L}^0(\Omega\times[-\pi,\pi]^d)$ such that $f_s\to f$ in measure.
    \end{itemize}
    Moreover, in such a case, $\{A_\bfn\}_\bfn\in\mathcal{G}_\Omega$ and $f$ is its unbounded GLT symbol.
\end{thm}

\begin{proof}First of all, consider a regular exhaustion $\{\Omega_t\}_t$ of $\Omega$.

    \noindent{\textbf{Step 1.}} We prove that $(i)\Rightarrow(ii)$. 
    
    Indeed, assume that there exists $\{A_\bfn\}_\bfn\in\mathcal{M}_\Omega$, such that $\{\{A_{\bfn,s}\}_\bfn\}_s\overset{a.c.s.}{\longrightarrow}\{A_\bfn\}_\bfn$.

    By continuity with respect to the a.c.s. convergence of the restriction operators (recall Remark \ref{rem_restr_btw_domains}), for every $t$, it follows that
    \begin{equation*}
        \mathcal{R}_{\Omega,\Omega_t}(\{\{A_{\bfn,s}\}_\bfn\}_s)\overset{a.c.s.}{\longrightarrow}\mathcal{R}_{\Omega,\Omega_t}(\{A_\bfn\}_\bfn),\qquad s\to+\infty.
    \end{equation*}
    On the other hand, by Theorem \ref{perm_spectr_info}, for every $s$ and every $t$, it holds 
    \begin{equation*}
        \mathcal{R}_{\Omega,\Omega_t}(\{A_{\bfn,s}\}_\bfn)\sim_{\mathrm{GLT}}^{\Omega_t}f_s|_{\Omega_t\times[-\pi,\pi]^d}.
    \end{equation*}
    By Theorem \ref{isometry_red_GLT}, any reduced GLT algebra is closed under a.c.s. convergence. Moreover, the a.c.s. convergence of a sequence of reduced GLT is equivalent to the convergence in measure of the corresponding canonical symbols. Therefore, for every $t$, we have $\mathcal{R}_{\Omega,\Omega_t}(\{A_\bfn\}_\bfn)\in\mathcal{G}_{\Omega_t}$ and there exists a function $g_t\in\mathcal{L}^0(\Omega_t\times[-\pi,\pi]^d)$, such that
    \begin{itemize}
        \item $f_s|_{\Omega_t\times[-\pi,\pi]^d}\to g_t$ in measure as $s\to +\infty$;
        \item $\mathcal{R}_{\Omega,\Omega_t}(\{A_\bfn\}_\bfn)\sim_{\mathrm{GLT}}^{\Omega_t}g_t$.
    \end{itemize}
    Since the restriction operators factor through intermediate domains and by the uniqueness of the reduced GLT symbol, for every $t<t'$, we have $g_{t'}|_{\Omega_{t}\times[-\pi,\pi]^d}=g_t$ almost everywhere in $\Omega_t\times[-\pi,\pi]^d$. Thus, we can define a measurable function $f:\Omega\times[-\pi,\pi]^d\to\mathbb{C}$ by declaring it to be equal to $g_t$ on $\Omega_t\times[-\pi,\pi]^d$, for every $t$. 
    Thanks to Remark \ref{rem_obs_ext_operator}, for every $t$, we can write 
    \begin{equation*}
        \{A_{\bfn}\}_\bfn=\mathcal{E}_{\Omega_t,\Omega}(\mathcal{R}_{\Omega,\Omega_t}(\{A_{\bfn}\}_\bfn))+\{S_{\bfn,t}\}_\bfn,
    \end{equation*}
    with $\mathrm{rank}(S_{\bfn,t})\leq 2(d_{\bfn}^{\Omega}-d_\bfn^{\Omega_t})$.
    
    Applying item (ii) of Lemma \ref{dimensions}, it follows that $\{\mathcal{R}_{\Omega,\Omega_t}(\{A_\bfn\}_\bfn)\}_t$ is a g.a.c.s. for $\{A_\bfn\}_\bfn$ made with compatible reduced GLT sequences. By Definition \ref{def_UGLT}, we conclude that $\{A_\bfn\}_\bfn\in\mathcal{G}_\Omega$ with canonical symbol $f$.
    
    It remains to prove that $f_s\to f$ in measure in $\Omega$. Indeed, for every $\delta>0$, there exists $t$ such that $\Ld(\Omega\setminus\Omega_t)\leq\delta$. As a consequence, for every $\varepsilon>0$, we have
    \begin{equation*}
        \mathcal{L}^{2d}(\{|f_s-f|>\varepsilon\})\leq (2\pi)^d\Ld(\Omega\setminus\Omega_t)+\mathcal{L}^{2d}(\{|(f_s-f)|_{\Omega_t}|>\varepsilon\})=(2\pi)^d\delta+\mathcal{L}^{2d}(\{|f_s|_{\Omega_t}-g_t|>\varepsilon\}),
    \end{equation*}
    so that passing to the $\limsup$, we have
    \begin{equation*}
        \limsup\limits_{s\to+\infty}\mathcal{L}^{2d}(\{|f_s-f|>\varepsilon\})\leq (2\pi)^d\delta+\limsup\limits_{s\to+\infty}\mathcal{L}^{2d}(\{|f_s|_{\Omega_t\times[-\pi,\pi]^d}-g_t|>\varepsilon\})=(2\pi)^d\delta.
    \end{equation*}
    By the arbitrariness of $\delta>0$, we conclude that $f_s\to f$ in measure.

    We divide the proof of $(ii)\Rightarrow(i)$ in two steps.\\
    \noindent{\textbf{Step 2.}} We find a candidate a.c.s. limit for $\{\{A_{\bfn,s}\}_\bfn\}_s$ in $\mathcal{M}_\Omega$, and we prove that it is an unbounded GLT over $\Omega$, with canonical symbol $f$.
    
    Indeed, assume that there exists $f\in\mathcal{L}^0(\Omega\times[-\pi,\pi]^d)$ such that $f_s\to f$ in measure, as $s\to+\infty$. Applying again Theorem \ref{isometry_red_GLT} to the sequences $\{\mathcal{R}_{\Omega,\Omega_t}(\{A_{\bfn,s}\}_\bfn)\}_s$ for $t$ fixed, since $f_s|_{\Omega_t}\to f|_{\Omega_t}$ in measure, it follows that, for every $t$ there is a sequence $\{B_{\bfn,t}\}_\bfn\in\mathcal{G}_{\Omega_t}$, such that
    \begin{itemize}
        \item $\mathcal{R}_{\Omega,\Omega_t}(\{\{A_{\bfn,s}\}_\bfn\}_s)\overset{a.c.s.}{\longrightarrow}\{B_{\bfn,t}\}_\bfn,$ as $s\to+\infty;$
        \item $\{B_{\bfn,t}\}_\bfn\sim_{\mathrm{GLT}}^{\Omega_t}f|_{\Omega_t\times[-\pi,\pi]^d}$.
    \end{itemize}
    Now, by a.c.s. continuity of the restriction operators, for every $t<t'$, it follows
    \begin{equation*}
        \mathcal{R}_{\Omega_{t'},\Omega_t}(\{B_{\bfn,t'}\}_\bfn)-\{B_{\bfn,t}\}_\bfn\sim_{\mathrm{GLT}}^{\Omega_t}0.
    \end{equation*}
    Applying the extension operator $\mathcal{E}_{\Omega_t,\Omega}$ to the previous identity and using the a.c.s. continuity, we get
    \begin{equation}\label{acs_eq_1}
        \mathcal{E}_{\Omega_t,\Omega}(\mathcal{R}_{\Omega_{t'},\Omega_t}(\{B_{\bfn,t'}\}_\bfn))-\mathcal{E}_{\Omega_t,\Omega}(\{B_{\bfn,t}\}_\bfn)\sim_{\sigma}0.
    \end{equation}
    As usual, by Remark \ref{rem_obs_ext_operator}, for every $t<t'$, we can write
    \begin{equation}\label{acs_eq_2}
        \{B_{\bfn,t'}\}=\mathcal{E}_{\Omega_{t'},\Omega}(\mathcal{R}_{\Omega_{t'},\Omega_t}(\{B_{\bfn,t'}\}_\bfn))+\{S_{\bfn,t,t'}\}_\bfn,
    \end{equation}
    with $\mathrm{rank}(S_{\bfn,t,t'})\leq 2(d_\bfn^{\Omega_{t'}}-d_\bfn^{\Omega})$.
    
    Since the extension operator $\mathcal{E}_{\Omega_{t'},\Omega}$ is rank preserving, by definition of $d_{a.c.s.}$ together with item (iii) of Lemma \ref{dimensions}, this implies that 
    \begin{align}\label{acs_eq_3}
        d_{a.c.s.}(\mathcal{E}_{\Omega_t,\Omega}(\{S_{\bfn,t,t'}\}_\bfn),\{0\}_\bfn)&=\limsup\limits_{\bfn\to\infty} p(\mathcal{E}_{\Omega_t,\Omega}(\{S_{\bfn,t,t'}\}_\bfn))\\
       \nonumber &\leq\limsup\limits_{\bfn\to\infty}2\frac{d_\bfn^{\Omega_{t'}}-d_\bfn^{\Omega_t}}{d_\bfn^{\Omega}}\\
       \nonumber &=2\frac{\Ld(\Omega_{t'})-\Ld(\Omega_t)}{\Ld(\Omega)}=2\frac{\Ld(\Omega_{t'}\setminus\Omega_t)}{\Ld(\Omega)}.
    \end{align}
    Combining now Equations \eqref{acs_eq_1}-\eqref{acs_eq_3}, for $t<t'$, we obtain
    \begin{align*}
        d_{a.c.s.}(\mathcal{E}_{\Omega_{t'},\Omega}(\{B_{\bfn,t'}\}_\bfn),\mathcal{E}_{\Omega_t,\Omega}(\{B_{\bfn,t}\}_\bfn))&\leq d_{a.c.s.}(\mathcal{E}_{\Omega_{t'},\Omega}(\{B_{\bfn,t'}\}_\bfn),\mathcal{E}_{\Omega_t,\Omega}(\mathcal{R}_{\Omega_{t'},\Omega_t}(\{B_{\bfn,t'}\}_\bfn)))\\
        &+d_{a.c.s.}(\mathcal{E}_{\Omega_t,\Omega}(\mathcal{R}_{\Omega_{t'},\Omega_t}(\{B_{\bfn,t'}\}_\bfn)),\mathcal{E}_{\Omega_t,\Omega}(\{B_{\bfn,t}\}_\bfn))\\
        &=d_{a.c.s.}(\mathcal{E}_{\Omega_t,\Omega}(\{S_{\bfn,t,t'}\}_\bfn),0)+0\\
        &\leq 2\frac{\Ld(\Omega_{t'}\setminus\Omega_t)}{\Ld(\Omega)}.
    \end{align*}
    In addition, since 
    \begin{equation*}
        \lim\limits_{(t,t')\to+\infty}2\frac{\Ld(\Omega_{t'}\setminus\Omega_t)}{\Ld(\Omega)}=0,
    \end{equation*}
   it follows that $\{\mathcal{E}_{\Omega_t,\Omega}(\{B_{\bfn,t}\}_\bfn)\}_t$ is a.c.s.-Cauchy in $\mathcal{M}_\Omega$. By Theorem \ref{completeness_acs_conv}, there exists $\{A_\bfn\}_\bfn\in\mathcal{M}_\Omega$, such that $\mathcal{E}_{\Omega_t,\Omega}(\{B_{\bfn,t}\}_\bfn)\overset{a.c.s.}{\longrightarrow}\{A_\bfn\}_\bfn$. Since this is equivalent to $\{\{B_{\bfn,t}\}_\bfn\}_t$ being a g.a.c.s. for $\{A_\bfn\}_\bfn$, it follows that $\{A_\bfn\}_\bfn\in\mathcal{G}_{\Omega}$, with canonical symbol $f$. 

   \noindent{\textbf{Step 3.}} We prove that $\{\{A_{\bfn,s}\}_\bfn\}_s\overset{a.c.s.}{\longrightarrow}\{A_\bfn\}_\bfn$, as $s\to+\infty$.
   
   Since $f_s\to f$ in measure in $\Omega\times[-\pi,\pi]^d$, clearly, for every $t$, $f_s|_{\Omega_t\times[-\pi,\pi]^d}\to f|_{\Omega_t\times[-\pi,\pi]^d}$ in measure in $\Omega_t\times[-\pi,\pi]^d$. Nevertheless, we need a more refined estimate of the convergence of the restrictions.
   
   Fix $s$ and $\varepsilon>0$ and, by definition of $d_{m,\Omega}$, there are two measurable sets $F_\varepsilon,H_\varepsilon\subset\Omega\times[-\pi,\pi]^d$, such that
   \begin{itemize}
       \item $F_\varepsilon\sqcup H_\varepsilon=\Omega\times[-\pi,\pi]^d$;
       \item $\frac{\mathcal{L}^{2d}(F_\varepsilon)}{(2\pi)^d\Ld(\Omega)}+\esssup\limits_{H_\varepsilon}|f_s-f|\leq d_{m,\Omega}(f_s,f)+\varepsilon$.
   \end{itemize}
   Now, we use their intersections with $\Omega_t\times[-\pi,\pi]^d$ to estimate from above the distance $d_{m,\Omega_t}$ between some suitable rescaling of the restrictions. Indeed,
   \begin{align*}
       d_{m,\Omega}(f_s,f)+\varepsilon&\geq \frac{\mathcal{L}^{2d}(F_\varepsilon)}{(2\pi)^d\Ld(\Omega)}+\esssup\limits_{H_\varepsilon}|f_s-f|\\
       &\geq \frac{\mathcal{L}^{2d}(F_\varepsilon\cap(\Omega_t\times[-\pi,\pi]^d))}{(2\pi)^d\Ld(\Omega)}+\esssup\limits_{H_\varepsilon\cap(\Omega_t\times[-\pi,\pi]^d)}|f_s-f|\\
       &=\frac{\Ld(\Omega_t)}{\Ld(\Omega)}\left(\frac{\mathcal{L}^{2d}(F_\varepsilon\cap(\Omega_t\times[-\pi,\pi]^d))}{(2\pi)^d\Ld(\Omega_t)}+\frac{\Ld(\Omega)}{\Ld(\Omega_t)}\cdot\esssup\limits_{H_\varepsilon\cap(\Omega_t\times[-\pi,\pi]^d)}|f_s-f|\right)\\
       &\geq \frac{\Ld(\Omega_t)}{\Ld(\Omega)}d_{m,\Omega_t}\left(\frac{\Ld(\Omega)}{\Ld(\Omega_t)}f_s|_{\Omega_t\times[-\pi,\pi]^d},\frac{\Ld(\Omega)}{\Ld(\Omega_t)}f|_{\Omega_t\times[-\pi,\pi]^d}\right).
   \end{align*}
   By the arbitrariness of $\varepsilon>0$, and reversing the inequality, we obtain
   \begin{equation}\label{acs_eq_4}
       d_{m,\Omega_t}\left(\frac{\Ld(\Omega)}{\Ld(\Omega_t)}f_s|_{\Omega_t\times[-\pi,\pi]^d},\frac{\Ld(\Omega)}{\Ld(\Omega_t)}f|_{\Omega_t\times[-\pi,\pi]^d}\right)\leq \frac{\Ld(\Omega)}{\Ld(\Omega_t)} d_{m,\Omega}(f_s,f).
   \end{equation}
   On the other hand, similarly as in the proof of Theorem \ref{isometry_red_GLT}, for fixed $\bfn$ and $t$ and for any $B\in\mathbb{C}^{d_\bfn^{\Omega}\times d_\bfn^{\Omega}}$, we have
   \begin{equation}\label{acs_eq_5}
       p(E_{\bfn,\Omega_t,\Omega}(B))=\frac{d_\bfn^{\Omega_t}}{d_\bfn^{\Omega}}p\left(\frac{d_\bfn^{\Omega}}{d_\bfn^{\Omega_t}}B\right).
   \end{equation}
   Using Equation \eqref{acs_eq_5}, item (iii) of Lemma \ref{dimensions} and passing to the $\limsup$, for any $\{B_\bfn\}_\bfn\in\mathcal{M}_{\Omega_t}$, we have
   \begin{equation}\label{acs_eq_6}
       d_{a.c.s.}\left(\mathcal{E}_{\Omega_t,\Omega}(\{B_\bfn\}_\bfn),\{0\}_\bfn\right)=\frac{\Ld(\Omega_t)}{\Ld(\Omega)}d_{a.c.s.}\left(\left\{\frac{d_\bfn^{\Omega}}{d_\bfn^{\Omega_t}}B_\bfn\right\}_\bfn,\{0\}_\bfn\right).
   \end{equation}
   Now, we want to estimate the quantity $d_{a.c.s.}(\{A_{\bfn,s}\}_\bfn,\{A_{\bfn}\}_\bfn)$. In order to do so, we decompose their difference into several pieces and estimate each one of them singularly. First of all, by Remark \ref{rem_obs_ext_operator}, for every $t$ and every $s$, we can write
   \begin{equation*}
       \{A_{\bfn,s}-A_\bfn\}_\bfn =\mathcal{E}_{\Omega_t,\Omega}(\mathcal{R}_{\Omega,\Omega_t}(\{A_{\bfn,s}-A_\bfn\}_\bfn))+\{S_{\bfn,s,t}\}_\bfn,
   \end{equation*}
   with $\mathrm{rank}(S_{\bfn,s,t})\leq 2(d_\bfn^{\Omega}-d_\bfn^{\Omega_t})$. Using the linearity of the operators, we obtain 
   \begin{align}\label{acs_eq_extra}
       \{A_{\bfn,s}-A_\bfn\}_\bfn &=\mathcal{E}_{\Omega_t,\Omega}(\mathcal{R}_{\Omega,\Omega_t}(\{A_{\bfn,s}-A_\bfn\}_\bfn))+\{S_{\bfn,s,t}\}_\bfn\\
       \nonumber&=\left(\mathcal{E}_{\Omega_t,\Omega}(\mathcal{R}_{\Omega,\Omega_t}(\{A_{\bfn,s}\}_\bfn))-\mathcal{E}_{\Omega_t,\Omega}(\{B_{\bfn,t}\}_\bfn)\right)\\
       \nonumber&+\left(\mathcal{E}_{\Omega_t,\Omega}(\{B_{\bfn,t}\}_\bfn)-\mathcal{E}_{\Omega_t,\Omega}(\mathcal{R}_{\Omega,\Omega_t}(\{A_{\bfn}\}_\bfn))\right)+\{S_{\bfn,s,t}\}_\bfn\\
   \end{align}
   Then, by the rank estimate and similarly as above, we have
   \begin{equation}\label{acs_eq_7}
       d_{a.c.s.}(\{S_{\bfn,s,t}\}_\bfn,\{0\}_\bfn)\leq 2\frac{\Ld(\Omega\setminus\Omega_t)}{\Ld(\Omega)}.
   \end{equation}
   Moreover, as a byproduct of the proof that $\{A_\bfn\}_\bfn$ is an unbounded GLT over $\Omega$, and since $\mathcal{R}_{\Omega,\Omega_t}(\{A_\bfn\}_\bfn)-\{B_{\bfn,t}\}_\bfn\sim_{\mathrm{GLT}}^{\Omega_t}0$ for every $t$,
   \begin{equation}\label{acs_eq_8}
       d_{a.c.s.}(\mathcal{E}_{\Omega_t,\Omega}(\{B_{\bfn,t}\}_\bfn),\mathcal{E}_{\Omega_t,\Omega}(\mathcal{R}_{\Omega,\Omega_t}(\{A_\bfn\}_\bfn)))=d_{a.c.s.}(\{B_{\bfn,t}\}_\bfn,\mathcal{R}_{\Omega,\Omega_t}(\{A_\bfn\}_\bfn))=0.
   \end{equation}
   Finally, we consider the following rescaled sequences 
   \begin{itemize}
       \item for every $s$, $\mathcal{R}_{\Omega,\Omega_t}\left(\left\{\frac{d_\bfn^{\Omega}}{d_\bfn^{\Omega_t}}A_{\bfn,s}\right\}_\bfn\right)\sim_{\mathrm{GLT}}^{\Omega_t}\frac{\Ld(\Omega)}{\Ld(\Omega_t)}f_s|_{\Omega_t\times[-\pi,\pi]^d}$;
       \item $\left\{\frac{d_\bfn^{\Omega}}{d_\bfn^{\Omega_t}}B_{\bfn,t}\right\}_\bfn\sim_{\mathrm{GLT}}^{\Omega_t}\frac{\Ld(\Omega)}{\Ld(\Omega_t)}f|_{\Omega_t\times[-\pi,\pi]^d}$.
   \end{itemize}
   Since these are reduced GLT, we can apply to them Theorem \ref{isometry_red_GLT}, so that the a.c.s. distance between any two of them is equal to the distance with respect to $d_{m,\Omega_t}$ of the corresponding symbols. 
   Therefore, using Equations \eqref{acs_eq_4} and \eqref{acs_eq_6}, we obtain 
   \begin{align}\label{acs_eq_9}
       \nonumber d_{a.c.s.}\left(\mathcal{E}_{\Omega_t,\Omega}(\mathcal{R}_{\Omega,\Omega_t}(\{A_{\bfn,s}\}_\bfn)),\mathcal{E}_{\Omega_t,\Omega}(\{B_{\bfn,t}\}_\bfn)\right)&=\frac{\Ld(\Omega_t)}{\Ld(\Omega)}d_{a.c.s.}\left(\mathcal{R}_{\Omega,\Omega_t}\left(\left\{\frac{d_\bfn^{\Omega}}{d_\bfn^{\Omega_t}}A_{\bfn,s}\right\}_\bfn\right),\left\{\frac{d_\bfn^{\Omega}}{d_\bfn^{\Omega_t}}B_{\bfn,t}\right\}_\bfn\right)\\
         &=\frac{\Ld(\Omega_t)}{\Ld(\Omega)}d_{m,\Omega_t}\left(\frac{\Ld(\Omega)}{\Ld(\Omega_t)}f_s|_{\Omega_t\times[-\pi,\pi]^d},\frac{\Ld(\Omega)}{\Ld(\Omega_t)}f|_{\Omega_t\times[-\pi,\pi]^d}\right)\\
        \nonumber &\leq \frac{\Ld(\Omega_t)}{\Ld(\Omega)}\frac{\Ld(\Omega)}{\Ld(\Omega_t)}d_{m,\Omega}(f_s,f)=d_{m,\Omega}(f_s,f).
   \end{align}
   To conclude, we compute the a.c.s. distance between $\{A_{\bfn,s}\}_\bfn$ and $\{A_\bfn\}_\bfn$, by using the expansion given by Equation \eqref{acs_eq_extra}, together with the estimates in Equations \eqref{acs_eq_7}-\eqref{acs_eq_9}. Indeed, for every $t$, we have 
   \begin{align}\label{acs_eq_10}
       d_{a.c.s.}\left(\{A_{\bfn,s}\}_\bfn, \{A_\bfn\}_\bfn\right)&\leq d_{a.c.s.}\left(\mathcal{E}_{\Omega_t,\Omega}(\mathcal{R}_{\Omega,\Omega_t}(\{A_{\bfn,s}\}_\bfn)),\mathcal{E}_{\Omega_t,\Omega}(\{B_{\bfn,t}\}_\bfn)\right)\\
       \nonumber&+d_{a.c.s.}(\mathcal{E}_{\Omega_t,\Omega}(\{B_{\bfn,t}\}_\bfn),\mathcal{E}_{\Omega_t,\Omega}(\mathcal{R}_{\Omega,\Omega_t}(\{A_\bfn\}_\bfn)))\\
       \nonumber&+d_{a.c.s.}(\{S_{\bfn,s,t}\}_\bfn,\{0\}_\bfn)\\
       \nonumber&\leq d_{m,\Omega}(f_s,f)+ 0 + 2\frac{\Ld(\Omega\setminus\Omega_t)}{\Ld(\Omega)}.
   \end{align}
   Since Equation \eqref{acs_eq_10} holds for every $t$, by letting $t\to+\infty$, we obtain
   \begin{equation*}
       d_{a.c.s.}\left(\{A_{\bfn,s}\}_\bfn, \{A_\bfn\}_\bfn\right)\leq d_{m,\Omega}(f_s,f),
   \end{equation*}
   which concludes the proof, since $d_{m,\Omega}(f_s,f)\to 0$ as $s\to+\infty$.
   \end{proof}

\end{document}
